\section{Analysis of Model-\texorpdfstring{$B_2$}{B2} (\texorpdfstring{$\gamma_1>0$}{gamma1>0}, \texorpdfstring{$\gamma_2=0$}{gamma2=0}, \texorpdfstring{$\theta_1=\theta_2=0$}{theta1=theta2=0})}\label{sec:model_b2}

Model-$B_2$ keeps one-way jockeying from class-$1$ to class-$2$ ($\gamma_1>0$, $\gamma_2=0$) and no abandonment ($\theta_1=\theta_2=0$). This model is more tractable than Model-$B$, as all the terms with derivatives in $y$ disappear, leaving us a first-order linear ODE in $x$, for each fixed $y\in(0,1]$. We will solve this using the integrating factor (\S\ref{sssec:if}). 

\textbf{Stability.} As in Models~$A$ and~$B$, $\rho = (\lambda_1+\lambda_2)/\mu < 1$. Jockeying conserves the total customer count and does not affect the stability condition.
Figure~\ref{fig:model-b2-queue} shows the queue layout.
\begin{figure}[H]
    \centering
    \resizebox{0.62\textwidth}{!}{%
        % Model B₂: γ₂=0, θ₁=θ₂=0, γ₁>0
% One-way jockeying: Class-1 can jockey into Class-2 queue only (dashed, downward).
% Class-2 customers cannot jockey back.
\begin{tikzpicture}[>=Latex, thick]

    % ── Queue 1 (top, Class-1, priority) ──────────────────────────────────────
    \draw (0, 1.9) -- (2, 1.9) -- (2, 1.1) -- (0, 1.1);
    \foreach \x in {0.2, 0.4, ..., 1.8} { \draw[thin] (\x, 1.8) -- (\x, 1.2); }
    \draw[->] (-1.5, 1.5) -- (-0.1, 1.5) node[midway, above] {$\lambda_1$};

    % ── Queue 2 (bottom, Class-2) ──────────────────────────────────────────────
    \draw (0, -1.1) -- (2, -1.1) -- (2, -1.9) -- (0, -1.9);
    \foreach \x in {0.2, 0.4, ..., 1.8} { \draw[thin] (\x, -1.2) -- (\x, -1.8); }
    \draw[->] (-1.5, -1.5) -- (-0.1, -1.5) node[midway, above] {$\lambda_2$};

    % ── One-way jockeying (1→2 only) ─────────────────────────────────────────
    \draw[->, dashed] (1.0, 1.0) -- (1.0, -1.0) node[midway, left] {$\gamma_1$};
    % No upward γ₂ arrow

    % ── Server ────────────────────────────────────────────────────────────────
    \node[draw, circle, minimum size=1.2cm] (server) at (5, 0) {$\mu$};
    \draw[->] (2.2,  1.5) -- (server);
    \draw[->] (2.2, -1.5) -- (server);
    \draw[->] (server) -- (7, 0);

\end{tikzpicture}
%
    }
    \caption{Queue layout for Model-$B_2$. Only the $1\to2$ jockeying direction
    is active (dashed, downward): class-1 customers may defect to Queue~2,
    but no customer can move in the reverse direction ($\gamma_2=0$).}
    \label{fig:model-b2-queue}
\end{figure}

\subsection{Reduced fundamental equation}\label{ssec:B2:reduction}

Setting $\gamma_2=0$ and $\theta_1=\theta_2=0$ in the general fundamental equation~\eqref{eq:gen:fundamental_equation} (Lemma~\ref{lem:gen:fundamental_equation}) annihilates the derivative terms in $y$, leaving:
\begin{align}
    &\bigl[(\lambda_1+\lambda_2+\mu)xy - \mu y - \lambda_1 x^2 y - \lambda_2 xy^2\bigr]P
    + \gamma_1 xy(x-y)\,\frac{\partial P}{\partial x} \nonumber \\
    &\qquad = \mu(x-y)\,P_y(y) - \mu x(1-y)\,\pi(0,0).
    \label{eq:B2:fundamental}
\end{align}
For each fixed $y\in(0,1]$ this is a first-order linear ODE in $x$; the entire analysis
of this section is built on it.

\begin{theorem}\label{thm:B2:PGF}
    Consider the two-class non-preemptive priority queueing system of Model-$B_2$ with
    per-customer jockeying rate $\gamma_1>0$ from Queue~1 to Queue~2 ($\gamma_2=0$,
    $\theta_1=\theta_2=0$). Under the stability condition $\rho=\lambda/\mu<1$,
    $\lambda=\lambda_1+\lambda_2$, the boundary generating function is
    \begin{equation}
        P_y(y) = \frac{\mu\,\pi(0,0)}{\mu-\lambda y}
        \cdot\frac{{}_1F_1\!\bigl(\alpha(y)+1;\,b^\star(y);\,-\lambda_1 y/\gamma_1\bigr)}
                  {{}_1F_1\!\bigl(\alpha(y);\,b^\star(y);\,-\lambda_1 y/\gamma_1\bigr)},
        \label{eq:B2:Py_thm}
    \end{equation}
    and the bivariate PGF, for $0\le x\le y$, is
    \begin{align}
        P(x,y) &= \frac{\mu\,e^{\lambda_1 x/\gamma_1}\,x^{-\alpha}(y-x)^{-\beta}}{\gamma_1 y}
        \Bigl[P_y(y)\,\mathcal{H}_1(x,y)
              + (1-y)\pi(0,0)\,\mathcal{H}_2(x,y)\Bigr],
        \label{eq:B2:Pxy}
    \end{align}
    where ${}_1F_1$ is Kummer's confluent hypergeometric function, the exponents are
    \begin{equation}
        \alpha(y) = \frac{\mu}{\gamma_1 y}, \qquad
        \beta(y)  = \frac{(1-y)(\lambda y-\mu)}{\gamma_1 y}, \qquad
        b^\star(y) = \alpha+\beta+1 = \frac{\mu+\lambda(1-y)+\gamma_1}{\gamma_1},
        \label{eq:B2:exponents_thm}
    \end{equation}
    the auxiliary integrals are
    \begin{align}
        \mathcal{H}_1(x,y) &= \int_0^x e^{-\lambda_1 t/\gamma_1}\,t^{\,\alpha-1}(y-t)^{\beta}\,dt,
        \label{eq:B2:I1}\\
        \mathcal{H}_2(x,y) &= \int_0^x e^{-\lambda_1 t/\gamma_1}\,t^{\,\alpha}(y-t)^{\beta-1}\,dt,
        \label{eq:B2:I2}
    \end{align}
    and the empty-state probability $\pi(0,0)=\rho(1-\rho)$ is given by Corollary~\ref{cor:B2:pi00} below.
    On $y<x\le1$ the same expression~\eqref{eq:B2:Pxy} holds with $\mathcal{H}_1,\mathcal{H}_2$
    replaced by their finite-part continuations; that the right-hand side is genuinely analytic
    across the interior point $x=y$ is the content of the smoothness condition established in
    Step~(ii) of the proof.
\end{theorem}

In the same fashion as in for Model-$B$, the result below is stated as a corollary in order for it to stay at the same level of hierarchy as that of other models.
\begin{corollary}\label{cor:B2:pi00}
    The empty-system probability $\pi_0$ (server idle, no one waiting) and the
    both-queues-empty probability $\pi(0,0)$ (server busy, $N_1=N_2=0$) of Model-$B_2$ are
    \begin{equation*}
        \pi_0 = 1-\rho, \qquad \pi(0,0) = \rho(1-\rho),
    \end{equation*}
    identical to the Model-$A$ baseline.
\end{corollary}
As in for Model-$B$, the proof follows trivially by setting $x=y=z$ in our fundamental equation (Equation~\eqref{eq:B2:fundamental}). The surviving equation is identical to the homologous in Model-$A$, and the result is proven in Corollary~\ref{cor:A:pi0}.

The proof of Theorem~\ref{thm:B2:PGF} consists of three steps: (i) obtaining the the standard form and the integration factor, (ii) determining $P_y(y)$ by analyticity in $x=y$, and (iii) recovering $P(x,y)$.

\begin{proof}[Proof of Theorem~\ref{thm:B2:PGF}]

% ─────────────────────────────────────────────────────────────────
\noindent\textit{Step~(i): Standard form and the integrating factor.}
\medskip

Fix $y\in(0,1]$, cancel a factor $y$ on both sides and divide the reduced fundamental equation~\eqref{eq:B2:fundamental} by the derivative term in $x$. The equation takes standard form $\tfrac{dP}{dx}+p(x;y)\,P=q(x;y)$ of \S\ref{sssec:if} (Equation~\eqref{eq:ode:standard}), with
\begin{align}
    \begin{split}
        p(x;y) &= \frac{(\lambda_1+\lambda_2+\mu)x - \mu - \lambda_1 x^2 - \lambda_2 xy}{\gamma_1 x(x-y)},
        \label{eq:B2:pxy}\\[4pt]
    \end{split} \\
    \begin{split}
        q(x;y) &= \frac{\mu\,P_y(y)}{\gamma_1 xy} - \frac{\mu(1-y)\,\pi(0,0)}{\gamma_1 y(x-y)}.
        \label{eq:B2:qxy}
    \end{split}
\end{align}
Write the numerator of $p$ as a polynomial in $x$,
\begin{equation*}
    N(x) = -\lambda_1 x^2 + \bigl[(\lambda_1+\lambda_2+\mu)-\lambda_2 y\bigr]x - \mu
         = -\lambda_1\bigl(x-x^*(y)\bigr)\bigl(x-x^+(y)\bigr),
\end{equation*}
where the roots $x^*(y),x^+(y)$ are those found for the kernel in Model-$A$ (Equation~\eqref{eq:A:xstar}), with , with $\lambda_1 x^* x^+ = \mu$ following from Vieta's formula. The Partial Fraction Decomposition of $p = N(x)/\bigl[\gamma_1 x(x-y)\bigr]$ read as rollows:
\begin{align*}
    A &= -\lambda_1, \nonumber\\
    B &= \frac{N(0)}{0-y} = \frac{-\mu}{-y} = \frac{\mu}{y}, \nonumber\\
    C &= \frac{N(y)}{y} = \frac{-(\lambda y-\mu)(y-1)}{y}
       = \frac{(1-y)(\lambda y-\mu)}{y},
\end{align*}
with $N(y)=-\lambda y^2+(\lambda+\mu)y-\mu=-(\lambda y-\mu)(y-1)$. Thus
\begin{equation}
    p(x;y) = \frac{1}{\gamma_1}\!\left[
        -\lambda_1 + \frac{\mu/y}{x} + \frac{(1-y)(\lambda y-\mu)/y}{x-y}
    \right].
    \label{eq:B2:pfrac}
\end{equation}

Since we will have a factor $(x-y)$ due to jockeying. we work on the region $0\leq x\leq y$, where $x-y\le0$; the complementary region $y<x\le1$ is recovered later. On this region, a term $\beta\ln|x-y|$ becomes simply $\beta\ln(y.-x)$, so integrating the partial fraction \eqref{eq:B2:pfrac} yields, up to a constant,
\begin{align*}
    \int^{x}\! p(t;y)\,dt
    &= \frac{1}{\gamma_1}\Bigl[-\lambda_1 x + \frac{\mu}{y}\ln x
       + \frac{(1-y)(\lambda y-\mu)}{y}\,\ln(y-x)\Bigr] \\
    &= -\frac{\lambda_1}{\gamma_1}\,x + \alpha(y)\ln x + \beta(y)\ln(y-x),
\end{align*}
with the exponents $\alpha(y)=\mu/(\gamma_1 y)>0$ and $\beta(y)=(1-y)(\lambda y-\mu)/(\gamma_1 y)$of~\eqref{eq:B2:exponents_thm}. Exponentiating, and dropping the multiplicative constant to which an integrating factor is indifferent, gives our integrating factor
\begin{equation*}
    \mu_I(x;y) = e^{-\lambda_1 x/\gamma_1}\,x^{\,\alpha(y)}\,(y-x)^{\,\beta(y)},
\end{equation*}
real-valued in our region because we take the base $(y-x)>0$ rather than $(x-y)$; its reciprocal $1/\mu_I$ is the homogeneous solution, playing here the role that $P_h$ played for Model-$B$ (\S\ref{sssec:vop}). It indeed solves $\mu_I'/\mu_I=p$:
\begin{align*}
    \frac{\mu_I'(x;y)}{\mu_I(x;y)}
    &= -\frac{\lambda_1}{\gamma_1} + \frac{\alpha}{x} + \beta\,\frac{d}{dx}\ln(y-x) \\
    &= -\frac{\lambda_1}{\gamma_1} + \frac{\alpha}{x} + \frac{\beta}{x-y}
     = p(x;y).
\end{align*}
The general solution of \S\ref{sssec:if} (Equation~\eqref{eq:ode:general}) now applies, and its arbitrary function is fixed immediately: since $\alpha>0$ we have $\mu_I\to0$ as $x\to0^{+}$, so the boundary contribution $P(0,y)\,\mu_I(0;y)$ vanishes and integrating $\tfrac{d}{dx}[P\,\mu_I]=q\,\mu_I$ over $(0,x]$ leaves
\begin{equation*}
    P(x,y)\,\mu_I(x;y) = \int_0^x q(t;y)\,\mu_I(t;y)\,dt.
\end{equation*}
Multiplying $q$ from~\eqref{eq:B2:qxy} by $\mu_I$ and applying the identity
$(y-t)^{\beta}/(t-y)=-(y-t)^{\beta-1}$ to the second term, the integrand separates into two pieces,
\begin{equation*}
    q(t;y)\,\mu_I(t;y) = \frac{\mu}{\gamma_1 y}\Bigl[
        P_y(y)\,e^{-\lambda_1 t/\gamma_1}\,t^{\,\alpha-1}(y-t)^{\beta}
        + (1-y)\pi(0,0)\,e^{-\lambda_1 t/\gamma_1}\,t^{\,\alpha}(y-t)^{\beta-1}
    \Bigr],
\end{equation*}
With the auxiliary integrals $\mathcal{H}_1,\mathcal{H}_2$ of~\eqref{eq:B2:I1}--\eqref{eq:B2:I2} and dividing through by $\mu_I(x;y)=e^{-\lambda_1 x/\gamma_1}x^{\alpha}(y-x)^{\beta}$, the solution on the principal region $0\le x\le y$ is
\begin{equation}
    P(x,y) = \frac{\mu\,e^{\lambda_1 x/\gamma_1}\,x^{-\alpha}(y-x)^{-\beta}}{\gamma_1 y}
    \Bigl[P_y(y)\,\mathcal{H}_1(x,y) + (1-y)\pi(0,0)\,\mathcal{H}_2(x,y)\Bigr],
    \label{eq:B2:Pxy_real}
\end{equation}
which still contains the unknown boundary function $P_y(y)$; Step~(ii) determines it.

% ─────────────────────────────────────────────────────────────────
\medskip\noindent\textit{Step~(ii): Determine $P_y(y)$ by analyticity at $x=y$.}
\medskip

Under stability —$\rho\leq1$— $\mu-\lambda y>0$ for $y\in(0,1)$, so the second exponent is negative:
\begin{equation}
    \beta(y) = \frac{(1-y)(\lambda y-\mu)}{\gamma_1 y} < 0,
    \qquad 0<y<1.
    \label{eq:B2:beta_negative}
\end{equation}
So the prefactor $(y-x)^{-\beta}=(y-x)^{|\beta|}$ in~\eqref{eq:B2:Pxy_real} \emph{vanishes} as $x\to y^{-}$, while $\mathcal{H}_1,\mathcal{H}_2$ may \emph{diverge}: the product is an indeterminate $0\cdot\infty$ form, whose resolution fixes $P_y(y)$.

On the diagonal the value is known outright. Setting $x=y$ in the fundamental
equation~\eqref{eq:B2:fundamental} drops the jockeying and boundary terms (each carries a
factor $x-y$) and leaves Model~$A$'s balance, so
\begin{equation}
    P(y,y) = \frac{\rho(1-\rho)}{1-\rho y},
    \label{eq:B2:regular_value}
\end{equation}
the diagonal value~\eqref{eq:A:Pzz}.

\smallskip\noindent\textit{Smoothness condition and $P_y(y)$.}
Pull the analytic, non-vanishing prefactor of~\eqref{eq:B2:Pxy_real} into
$\Psi(x)=\mu\,e^{\lambda_1 x/\gamma_1}x^{-\alpha}/(\gamma_1 y)$, so that
\begin{equation*}
    P(x,y)=\Psi(x)\,(y-x)^{|\beta|}
        \bigl[\,P_y(y)\,\mathcal{H}_1(x,y)+(1-y)\pi(0,0)\,\mathcal{H}_2(x,y)\,\bigr],
    \qquad |\beta|=-\beta>0 .
\end{equation*}
We need only the leading behaviour of the two integrals as $x\to y^{-}$, not a full expansion.
Each splits into its endpoint value $\mathcal{H}_i(y,y)=\int_0^y(\cdots)\,dt$ and a remainder
set by the integrand near $t=y$ --- $\mathcal{H}_1$ carries $(y-t)^{\beta}$, $\mathcal{H}_2$
carries $(y-t)^{\beta-1}$ --- so that, with constants $c_1,c_2$ whose values we never need,
\begin{equation*}
    \mathcal{H}_1(x,y)=\mathcal{H}_1(y,y)+c_1\,(y-x)^{\beta+1}+\cdots, \qquad
    \mathcal{H}_2(x,y)=\mathcal{H}_2(y,y)+c_2\,(y-x)^{\beta}+\cdots .
\end{equation*}
Multiplied by $(y-x)^{-\beta}$, the $c_1,c_2$ remainders turn into a smooth $O(y-x)$ term and the
constant regular value~\eqref{eq:B2:regular_value}; only the endpoint values retain the bare,
non-integer power $(y-x)^{|\beta|}$ ($|\beta|\notin\mathbb{Z}$ generically). Since $P(\cdot,y)$
is a power series in $x$ --- hence smooth at the interior point $x=y<1$, whereas that power is
not --- its coefficient must vanish:
\begin{equation}
    P_y(y)\,\mathcal{H}_1(y,y) + (1-y)\,\pi(0,0)\,\mathcal{H}_2(y,y) = 0 ,
    \label{eq:B2:nobranch}
\end{equation}
the endpoint integrals being read, since $\beta<0$, as the Euler--Beta continuations found next.
To evaluate these endpoint integrals, substitute $t=ys$ ($dt=y\,ds$)
in~\eqref{eq:B2:I1}--\eqref{eq:B2:I2}, which rescales the two integrands and pulls out a common
$y^{\alpha+\beta}$,
\begin{align*}
    t^{\alpha-1}(y-t)^{\beta}\,dt &= y^{\alpha+\beta}\,s^{\alpha-1}(1-s)^{\beta}\,ds, \\
    t^{\alpha}(y-t)^{\beta-1}\,dt &= y^{\alpha+\beta}\,s^{\alpha}(1-s)^{\beta-1}\,ds,
\end{align*}
so that, with $c:=\lambda_1 y/\gamma_1$,
\[
    \mathcal{H}_1(y,y)=y^{\alpha+\beta}\!\int_0^1\! s^{\alpha-1}(1-s)^{\beta}e^{-cs}\,ds,
    \qquad
    \mathcal{H}_2(y,y)=y^{\alpha+\beta}\!\int_0^1\! s^{\alpha}(1-s)^{\beta-1}e^{-cs}\,ds.
\]
These are Euler integrals~\eqref{eq:kummer:euler},
$\int_0^1 t^{a-1}(1-t)^{b-1}e^{zt}\,dt=B(a,b)\,{}_1F_1(a;a+b;z)$, with
$(a,b,z)=(\alpha,\beta+1,-c)$ for $\mathcal{H}_1$ and $(\alpha+1,\beta,-c)$ for
$\mathcal{H}_2$, so $a+b=b^\star(y)$ in both. As $\beta<0$ by~\eqref{eq:B2:beta_negative}, the
right-hand side is read as its analytic continuation in $b$ beyond the convergence range $b>0$,
giving
\begin{align}
    \mathcal{H}_1(y,y)
        &= B(\alpha,\,\beta+1)\,y^{\alpha+\beta}\,
           {}_1F_1\!\bigl(\alpha;\,b^\star(y);\,-\lambda_1 y/\gamma_1\bigr),
    \label{eq:B2:fp1}\\[3pt]
    \mathcal{H}_2(y,y)
        &= B(\alpha+1,\,\beta)\,y^{\alpha+\beta}\,
           {}_1F_1\!\bigl(\alpha+1;\,b^\star(y);\,-\lambda_1 y/\gamma_1\bigr).
    \label{eq:B2:fp2}
\end{align}
Solving the smoothness condition~\eqref{eq:B2:nobranch} for $P_y(y)$ and inserting the
endpoint values~\eqref{eq:B2:fp1}--\eqref{eq:B2:fp2}, the common factor $y^{\alpha+\beta}$
cancels and
\begin{equation}
    P_y(y)
    = -(1-y)\,\pi(0,0)\,\frac{\mathcal{H}_2(y,y)}{\mathcal{H}_1(y,y)}
    = -(1-y)\,\pi(0,0)\,\frac{B(\alpha+1,\,\beta)}{B(\alpha,\,\beta+1)}\,
        \frac{{}_1F_1\!\bigl(\alpha+1;\,b^\star;\,-\lambda_1 y/\gamma_1\bigr)}
             {{}_1F_1\!\bigl(\alpha;\,b^\star;\,-\lambda_1 y/\gamma_1\bigr)}.
    \label{eq:B2:Py_ratio}
\end{equation}
The Beta ratio reduces, via $\Gamma(\alpha+1)=\alpha\Gamma(\alpha)$ and
$\Gamma(\beta+1)=\beta\Gamma(\beta)$, to
\begin{align}
    \frac{B(\alpha+1,\,\beta)}{B(\alpha,\,\beta+1)}
    &= \frac{\Gamma(\alpha+1)\Gamma(\beta)}{\Gamma(\alpha+\beta+1)}
       \cdot\frac{\Gamma(\alpha+\beta+1)}{\Gamma(\alpha)\Gamma(\beta+1)} \nonumber\\
    &= \frac{\Gamma(\alpha+1)}{\Gamma(\alpha)}\cdot\frac{\Gamma(\beta)}{\Gamma(\beta+1)}
       = \frac{\alpha}{\beta},
    \label{eq:B2:beta_ratio}
\end{align}
and the surviving prefactor simplifies, using
$\alpha/\beta=\mu/\bigl[(1-y)(\lambda y-\mu)\bigr]$, as
\begin{align}
    -(1-y)\,\frac{\alpha}{\beta}
    &= -(1-y)\cdot\frac{\mu}{(1-y)(\lambda y-\mu)} \nonumber\\
    &= -\frac{\mu}{\lambda y-\mu} = \frac{\mu}{\mu-\lambda y}.
    \label{eq:B2:prefactor_combo}
\end{align}
Substituting~\eqref{eq:B2:beta_ratio}--\eqref{eq:B2:prefactor_combo}
into~\eqref{eq:B2:Py_ratio} gives the boundary generating function
\begin{align}
    P_y(y) &= \frac{\mu\,\pi(0,0)}{\mu-\lambda y}
    \cdot\frac{{}_1F_1\!\bigl(\alpha(y)+1;\,b^\star(y);\,-\lambda_1 y/\gamma_1\bigr)}
              {{}_1F_1\!\bigl(\alpha(y);\,b^\star(y);\,-\lambda_1 y/\gamma_1\bigr)},
    \label{eq:B2:Py}
\end{align}
confirming~\eqref{eq:B2:Py_thm}.  The prefactor equals the regular value $P(y,y)$
of~\eqref{eq:B2:regular_value}, and the Kummer ratio is the boundary correction.
As $y\to0^+$ the argument $-\lambda_1 y/\gamma_1\to0$, so the ratio tends to $1$ and
$P_y(0)=\pi(0,0)$, consistent with
$P_y(0)=\sum_{n_2}\pi(0,n_2)\cdot0^{n_2}=\pi(0,0)$.

% ─────────────────────────────────────────────────────────────────
\medskip\noindent\textit{Step~(iii): Recovery of $P(x,y)$.}
\medskip

Substituting~\eqref{eq:B2:Py} into the real-form expression~\eqref{eq:B2:Pxy_real}
gives, for $0\le x\le y$,
\begin{align*}
    P(x,y) &= \frac{\mu\,e^{\lambda_1 x/\gamma_1}\,x^{-\alpha}(y-x)^{-\beta}}{\gamma_1 y}
    \biggl[
        \frac{\mu\,\pi(0,0)}{\mu-\lambda y}
        \cdot\frac{{}_1F_1\!\bigl(\alpha+1;\,b^\star;\,-\tfrac{\lambda_1 y}{\gamma_1}\bigr)}
                  {{}_1F_1\!\bigl(\alpha;\,b^\star;\,-\tfrac{\lambda_1 y}{\gamma_1}\bigr)}
        \cdot\mathcal{H}_1(x,y)
    \nonumber\\[4pt]
    &\hspace{5.8cm}
        + (1-y)\pi(0,0)\cdot\mathcal{H}_2(x,y)
    \biggr],
\end{align*}
which is the statement of the theorem (Equation~\eqref{eq:B2:Pxy}).
The smoothness condition~\eqref{eq:B2:nobranch} guarantees that this extends
analytically across $x=y$, with diagonal value
$P(y,y)=\rho(1-\rho)/(1-\rho y)$ as computed in~\eqref{eq:B2:regular_value}.
On $y<x\le1$ the same formula holds with $\mathcal{H}_1,\mathcal{H}_2$ read as their
finite-part continuations; $\pi(0,0)=\rho(1-\rho)$ from Corollary~\ref{cor:B2:pi00} makes
the solution fully explicit.
The recovery of Model-$A$ in the limit $\gamma_1\to0^{+}$ is deferred to
Remark~\ref{rem:B2:singular_limit} below, as the limit is singular and merits
separate discussion.
\end{proof}

\begin{remark}[The singular limit $\gamma_1\to0^{+}$]\label{rem:B2:singular_limit}
The no-jockeying limit $\gamma_1\to0^{+}$ is a \emph{singular} perturbation: the coefficient
$\gamma_1 xy(x-y)$ of the highest derivative $\partial_x P$ vanishes, so one must check that the
convection term itself, not merely its coefficient, disappears. It does. Each $P$ is a
probability generating function, so $|P|\le1$ on the closed unit polydisc uniformly in
$\gamma_1$, and Cauchy's estimate then bounds $\partial_x P$ uniformly on every compact subset
of the open polydisc; hence $\gamma_1 xy(x-y)\,\partial_x P\to0$ there. The fundamental
equation~\eqref{eq:B2:fundamental} therefore degenerates to the algebraic Model-$A$
equation~\eqref{eq:A:fundamental}, and the uniformly bounded family $\{P\}_{\gamma_1}$, being
normal, converges (in full, by uniqueness of the analytic solution) to the Model-$A$
PGF~\eqref{eq:A:PGF}. In particular the boundary function reduces to the Model-$A$
value~\eqref{eq:A:Py},
\begin{equation}
    P_y(y)\;\xrightarrow[\ \gamma_1\to0^{+}\ ]{}\;
    \pi(0,0)\,\frac{x^*(y)\,(1-y)}{x^*(y)-y},
    \label{eq:B2:Py_limit}
\end{equation}
with $x^*(y)$ the Model-$A$ kernel root~\eqref{eq:A:xstar}; this needs no special-function
asymptotics.

The same limit can be read off the closed form~\eqref{eq:B2:Py}. Its prefactor
$\mu\,\pi(0,0)/(\mu-\lambda y)=\pi(0,0)/(1-\rho y)$ is $\gamma_1$-independent,
so~\eqref{eq:B2:Py_limit} is equivalent to the Kummer ratio tending to
\begin{equation}
    \frac{{}_1F_1\!\bigl(\alpha+1;\,b^\star;\,-\lambda_1 y/\gamma_1\bigr)}
         {{}_1F_1\!\bigl(\alpha;\,b^\star;\,-\lambda_1 y/\gamma_1\bigr)}
    \;\xrightarrow[\ \gamma_1\to0^{+}\ ]{}\;
    (1-\rho y)\,\frac{x^*(y)\,(1-y)}{x^*(y)-y}.
    \label{eq:B2:ratio_limit}
\end{equation}
The factor $(1-\rho y)$ shows the ratio does \emph{not} tend to the bare
$x^*(y)(1-y)/(x^*(y)-y)$; the discrepancy is exactly the prefactor. This matches the
first-parameter three-term recurrence for ${}_1F_1$~\cite[\S13.3]{dlmf}: as $\gamma_1\to0^{+}$
the parameters $\alpha$, $b^\star$ and $c=\lambda_1 y/\gamma_1$ all grow like $1/\gamma_1$, and
the recurrence reduces at leading order to a quadratic whose admissible root
is~\eqref{eq:B2:ratio_limit}.
\end{remark}

% ─────────────────────────────────────────────────────────────────
\subsection{Mean queue lengths and mean waiting times}

\subsubsection*{Total queue length from the diagonal}

Setting $x=y=z$ in the reduced fundamental equation~\eqref{eq:B2:fundamental}, the
jockeying convection term $\gamma_1 xy(x-y)\partial_xP$ vanishes at $x=y=z$
(factor $(x-y)=0$), as does the boundary term $\mu(x-y)P_y(y)$.  The remaining
balance is identical to Model-$A$'s, so the same diagonal cancellation gives
\[
    P(z,z) = \frac{\rho(1-\rho)}{1-\rho z},
\]
the same formula as~\eqref{eq:A:Pzz}.  Since $P(z,z)=\mathbb{E}[z^{N_1+N_2}]=\mathbb{E}[z^{N}]$
is the PGF of the total in-queue count $N$, the generating-function identity
$\tfrac{d}{dz}P(z,z)\big|_{z=1}=\mathbb{E}[N_1]+\mathbb{E}[N_2]$ holds;
differentiating by the quotient rule and evaluating at $z=1$:
\begin{align*}
    \mathbb{E}[N_1]+\mathbb{E}[N_2]
    &= \left.\frac{d}{dz}\,\frac{\rho(1-\rho)}{1-\rho z}\right|_{z=1}
     = \frac{\rho(1-\rho)\,\rho}{(1-\rho z)^2}\bigg|_{z=1} \\
    &= \frac{\rho^2(1-\rho)}{(1-\rho)^2}
     = \frac{\rho^2}{1-\rho}.
\end{align*}
One-way jockeying does not alter the total mean queue length relative to
Model-$A$, since the diagonal PGF is unchanged.

\subsubsection*{Class-1 mean queue length}

We extract $\mathbb{E}[N_1]=\tfrac{\partial P}{\partial x}\big|_{(1,1)}$ by first
reducing the fundamental equation~\eqref{eq:B2:fundamental} at $y=1$, and then
differentiating the resulting marginal PGF at $x=1$.

\medskip\noindent\textit{Step~1: marginal PGF $P(x,1)$.}\enspace
Evaluating~\eqref{eq:B2:fundamental} at $y=1$, the bracket multiplying $P$ becomes
$(\lambda_1+\mu)x-\mu-\lambda_1 x^2=-\lambda_1(x-1)(x-1/\rho_1)$, the convection
coefficient is $\gamma_1 x(x-1)$, and the right-hand side reduces to $\mu(x-1)P_y(1)$
(the $\pi(0,0)$ term carries the factor $(1-y)=0$). Every term therefore shares the common
factor $(x-1)$; cancelling it leaves the first-order linear ODE in $x$:
\[
    -\lambda_1(x-1/\rho_1)\,P(x,1) + \gamma_1 x\,\frac{d}{dx}P(x,1) = \mu\,P_y(1).
\]
The integrating factor for this ODE is $e^{-\lambda_1 x/\gamma_1}\,x^{\,\mu/\gamma_1}$
(noting $\lambda_1/\rho_1=\mu$), and integrating from $0$ to $x$ (the boundary term
vanishes since $0^{\mu/\gamma_1}=0$):
\[
    P(x,1)\,e^{-\lambda_1 x/\gamma_1}\,x^{\,\mu/\gamma_1}
    = \frac{\mu\,P_y(1)}{\gamma_1}
    \int_0^x e^{-\lambda_1 t/\gamma_1}\,t^{\,\mu/\gamma_1-1}\,dt.
\]
Rearranging and writing $\alpha=\mu/\gamma_1$ and
$\mathcal{H}(x)=\int_0^x e^{-\lambda_1 t/\gamma_1}\,t^{\,\alpha-1}\,dt$:
\[
    P(x,1) = \frac{\mu\,P_y(1)}{\gamma_1}\,e^{\lambda_1 x/\gamma_1}\,x^{-\alpha}\,\mathcal{H}(x).
\]
The constant $P_y(1)$ follows from $P(1,1)=1-\pi_0=\rho$, which gives
$P_y(1)=\rho\gamma_1/\bigl(\mu\,e^{\lambda_1/\gamma_1}\,\mathcal{H}(1)\bigr)$.

\medskip\noindent\textit{Step~2: differentiate $P(x,1)$ at $x=1$.}\enspace
Writing $P(x,1)=C\cdot f(x)\cdot\mathcal{H}(x)$ with $C=\mu P_y(1)/\gamma_1$ and
$f(x)=e^{\lambda_1 x/\gamma_1}\,x^{-\alpha}$, the product rule together with
Leibniz's rule ($d\mathcal{H}/dx = e^{-\lambda_1 x/\gamma_1}\,x^{\alpha-1}$, the integrand
at the upper limit) gives
\[
    \frac{d}{dx}P(x,1)
    = C\Bigl[f'(x)\,\mathcal{H}(x)+f(x)\,e^{-\lambda_1 x/\gamma_1}\,x^{\alpha-1}\Bigr].
\]
The logarithmic derivative $f'(x)/f(x)=\lambda_1/\gamma_1-\alpha/x=(\lambda_1 x-\mu)/(\gamma_1 x)$,
and the second product term simplifies as
$f(x)\,e^{-\lambda_1 x/\gamma_1}x^{\alpha-1}=x^{-1}$, so
\[
    \frac{d}{dx}P(x,1)
    = \frac{\mu\,P_y(1)}{\gamma_1}\left[
        e^{\lambda_1 x/\gamma_1}\,x^{-\alpha}\,\frac{\lambda_1 x-\mu}{\gamma_1 x}\,\mathcal{H}(x)
        + x^{-1}
    \right].
\]

\medskip\noindent\textit{Step~3: evaluate at $x=1$.}\enspace
\begin{align}
    \mathbb{E}[N_1]
    &= \left.\frac{d}{dx}P(x,1)\right|_{x=1} \nonumber\\
    &= \frac{\mu\,P_y(1)}{\gamma_1}\left[
        e^{\lambda_1/\gamma_1}\,\frac{\lambda_1-\mu}{\gamma_1}\,\mathcal{H}(1) + 1
    \right],
    \label{eq:B2:EN1}
\end{align}
where $\mathcal{H}(1)=\int_0^1 e^{-\lambda_1 t/\gamma_1}\,t^{\,\mu/\gamma_1-1}\,dt
=(\gamma_1/\lambda_1)^{\mu/\gamma_1}\,\gamma(\mu/\gamma_1,\,\lambda_1/\gamma_1)$
and $\gamma(\cdot,\cdot)$ is the lower incomplete gamma function.  The
factor $(\lambda_1-\mu)<0$ (since $\rho_1=\lambda_1/\mu<1$), confirming that
$e^{\lambda_1/\gamma_1}(\lambda_1-\mu)\mathcal{H}(1)/\gamma_1 + 1 > 0$.
No elementary simplification of~\eqref{eq:B2:EN1} is known; the expression
is evaluated numerically.  As $\gamma_1\to0^+$ the
Model-$A$ value~\eqref{eq:A:EN1} is recovered.

\subsubsection*{Class-2 mean queue length by subtraction}

\[
    \mathbb{E}[N_2]
    = \bigl(\mathbb{E}[N_1]+\mathbb{E}[N_2]\bigr)-\mathbb{E}[N_1]
    = \frac{\rho^2}{1-\rho}-\mathbb{E}[N_1].
\]
Since jockeying defections reduce the class-1 queue,
$\mathbb{E}[N_1]$ is smaller than the Model-$A$ value $\rho\rho_1/(1-\rho_1)$,
so $\mathbb{E}[N_2]$ correspondingly exceeds the Model-$A$ value.

\subsubsection*{Mean waiting times via Little's Law}

\textbf{Little's Law} applies here without qualification. It is a rate-conservation
identity --- for any subsystem in which customers are neither created nor destroyed
internally, the time-average number present equals the long-run entry rate times the mean
time an entrant spends inside --- and it holds for \emph{any} positive-recurrent system with
finite mean queue length (here ensured by $\rho<1$, which gives
$\mathbb{E}[N_1]\le\mathbb{E}[N]=\rho^2/(1-\rho)<\infty$ and so makes the jockeying rate
$\gamma_1\mathbb{E}[N_1]$ below a genuine finite rate), requiring neither Poisson input nor
independence. In particular it remains valid even
though one-way jockeying has destroyed the per-class Poisson/renewal structure (the very
obstruction that ruled out a probabilistic route in Step~(ii)). We apply it to each queue
(the waiting line, excluding the customer in service), with $\mathbb{E}[W_i]$ understood as
the mean time a customer spends in Queue~$i$ averaged over all customers that enter it ---
an arrival routed directly into service contributing a zero-length visit.

\emph{Class~1.} Every class-1 arrival is counted as an entrant to Queue~1 at rate
$\lambda_1$ (those finding the server idle enter service at once, contributing zero waiting
time), and each leaves Queue~1 either by entering service or by jockeying to Queue~2. Since
no other stream feeds Queue~1, its throughput is exactly $\lambda_1$, and Little's Law gives
\begin{equation}
    \mathbb{E}[W_1] = \frac{\mathbb{E}[N_1]}{\lambda_1}.
    \label{eq:B2:EW}
\end{equation}

\emph{Class~2.} Queue~2 has \emph{two} input streams: the class-2 Poisson arrivals at rate
$\lambda_2$ and the jockeying transfers from Queue~1. Because each of the $N_1$ waiting
class-1 customers jockeys independently at rate $\gamma_1$, the long-run transfer rate is
$\gamma_1\,\mathbb{E}[N_1]$ (a genuine rate by the ergodic theorem, equal to the time
average of the instantaneous jockeying intensity $\gamma_1 N_1$). The throughput of Queue~2
is therefore $\lambda_2+\gamma_1\,\mathbb{E}[N_1]$, and Little's Law reads
$\mathbb{E}[N_2]=\bigl(\lambda_2+\gamma_1\,\mathbb{E}[N_1]\bigr)\,\mathbb{E}[W_2]$, so that
\begin{equation}
    \mathbb{E}[W_2] = \frac{\mathbb{E}[N_2]}{\lambda_2+\gamma_1\,\mathbb{E}[N_1]}.
    \label{eq:B2:EW2}
\end{equation}
Here $\mathbb{E}[W_2]$ is the mean time in the class-2 line averaged over \emph{all} Queue-2
entrants of both streams; a jockeyed customer's Queue-2 clock starts at the transfer
instant, so its total wait is an initial Queue-1 segment (of mean $\mathbb{E}[W_1]$) followed
by this Queue-2 segment. Thus~\eqref{eq:B2:EW2} is the per-Queue-2-entrant residence time,
not the end-to-end delay of a tagged class-2 \emph{arrival}.

The jockeying inflow is what distinguishes~\eqref{eq:B2:EW2} from the single-stream
Models~$A$ and~$C_2$: the naive denominator $\lambda_2$ would undercount the entrants to
Queue~2 and is correct only in the no-jockeying limit. Indeed, as $\gamma_1\to0^{+}$ the
transfer rate $\gamma_1\,\mathbb{E}[N_1]\to0$, so~\eqref{eq:B2:EW2} reduces to
$\mathbb{E}[N_2]/\lambda_2$ and both~\eqref{eq:B2:EW} and~\eqref{eq:B2:EW2} recover the
Model-$A$ values~\eqref{eq:A:EW1}--\eqref{eq:A:EW2}.
